\begin{textsample}{NEG Dim 1 – 2015\_c5 – Score 6.00 – 2015\_c5/enade\_2015\_tec\_des\_inter\_2e210f18ec.txt}  \label{ex:f1_neg_344}
CAPÍTULO 5 DISTRIBUIÇÃO DOS CONCEITOS 5.1 PANORAMA NACIONAL DA DISTRIBUIÇÃO DOS CONCEITOS O Conceito Enade é um indicador de qualidade que avalia o desempenho dos estudantes a partir dos resultados do Enade.
Os valores possíveis vão de 1 ( pior situação ) a 5 ( melhor situação ).
As Unidades de Observação com apenas um ou sem nenhum concluinte participante não obtêm o Conceito Enade, ficando Sem Conceito ( SC ).
A Tabela 5.1 apresenta a quantidade e distribuição de cursos de Tecnologia em Design de Interiores participantes do Enade/2015, por faixa de conceito e Grande Região.
Mais uma vez, ressaltando que a diferença entre os cursos tabulados neste capítulo e no capítulo 2 corresponde aos cursos sem conceito, em princípio, aqueles sem alunos concluintes que participassem da prova ou que tivessem tão somente um aluno concluinte18.
Observando -se os dados da Tabela 5.1, nota -se que, dos 82 cursos participantes, 29 ( 35,4 \% ) classificaram -se com conceito 3, o valor modal.
Este foi também o conceito modal nas regiões Nordeste ( 50,0 \% ), Sudeste ( 36,8 \% ) e na região Norte ( 50,0 \% ), onde os dois cursos foram alocados a conceitos contíguos.
Nas regiões Sul e Centro-Oeste, a moda foi o conceito 4 ( 45,0 \% e 37,5 \%, respectivamente ).
O conceito 4 foi o segundo mais frequente em nível nacional ( 28,0 \%, correspondendo a 23 cursos ) e o conceito 2, o terceiro ( 19,5 \%, correspondendo a 16 cursos ).
Houve, ainda, dez cursos ( 12,2 \% ) que receberam conceito 5 e três ( 3,7 \% ) cursos que receberam conceito 1.
Dos 82 cursos de Tecnologia em Design de Interiores, apenas um curso ficou sem conceito ( SC ). 18 Estes cursos com somente um concluinte não foram considerados no capítulo 2.
A região Norte participou com dois cursos ( 2,4 \% do total nacional ).
Nesta região, dois cursos foram avaliados com conceitos 2 e 3 ( um curso, \textbf{cada} ), correspondendo a 50,0 \%, \textbf{cada}.
Nenhum curso ficou sem conceito ( SC ) ou recebeu os demais conceitos nessa região.
A região Nordeste participou com 14 cursos ( 17,1 \% do total nacional ).
Nesta região, sete cursos, 50,0 \% em termos regionais, obtiveram conceito 3, o conceito modal para a região.
Os conceitos 4 e 5 foram atribuídos a, respectivamente, três e quatro cursos ( 21,4 \% e 28,6 \% ).
Nenhum curso ficou sem conceito ( SC ) ou recebeu os conceitos 1 ou 2, nessa região.
Dos 38 cursos participantes da região Sudeste ( 46,3 \% do total nacional ), o conceito 3 também foi o conceito modal, tendo sido obtido por 14 ( 36,8 \% ) cursos.
O conceito 1 foi atribuído a dois cursos ( 5,3 \% ) e os conceitos 2 e 4, a dez e oito cursos, correspondendo a 26,3 \% e 21,1 \%, respectivamente.
O conceito 5 foi atribuído a três cursos ( 7,9 \% ).
E um curso ( 2,6 \% ) ficou sem conceito ( SC ).
Os 20 cursos da região Sul ficaram distribuídos em quase toda a faixa de conceitos.
Há predominância do conceito 4 atribuído a 45,0 \% dos cursos ( conceito modal ), ou seja, a nove dos 20 cursos participantes na região Sul.
O conceito 2 foi atribuído a três cursos ( 15,0 \% ) e o conceito 3, a seis cursos ( 30,0 \% ).
O conceito 5 foi recebido por dois cursos ( 10,0 \% ).
Nenhum curso na região Sul ficou sem conceito ( SC ) ou recebeu conceito 1.
Já dos oito cursos participantes na região Centro-Oeste, três receberam conceito 4, o conceito modal ( 37,5 \% ).
Nesta região, nenhum curso ficou sem conceito ( SC ).
Dois cursos ( 25,0 \% ) foram alocados no conceito 2 e três cursos foram avaliados com conceitos 1, 3 ou 5 ( um curso \textbf{cada} ), nessa região.
O Gráfico 5.1 apresenta a distribuição do Conceito Enade segundo Grande Região : poligonais mais à esquerda representam Grandes Regiões com pior distribuição, e poligonais mais à direita, Grandes Regiões com melhores conceitos.
A região Norte ( linha azul ) apresenta os cursos em IES com conceitos mais baixos, ao passo que a região Nordeste ( linha verde ) apresenta a distribuição com valores maiores.
Cursos em IES das regiões Sudeste ( linha vermelha ), Sul ( linha roxa ) e Centro-Oeste ( linha preta ) apresentam uma situação intermediária. 5.2 CONCEITOS POR CATEGORIA ADMINISTRATIVA E POR GRANDE REGIÃO A Tabela 5.2 apresenta a distribuição dos cursos participantes do Enade/2015 de Tecnologia em Design de Interiores, por Categoria Administrativa, de acordo com os conceitos por eles alcançados, segundo as Grandes Regiões brasileiras.
Dos 82 cursos participantes, apenas três ( 3,7 \% ) eram ministrados em Instituições Públicas e 79 ( 96,3 \% ), em Instituições Privadas.
De acordo com as informações da Tabela 5.2, em termos nacionais, entre Instituições Públicas e Privadas observa -se que, dos dez cursos que receberam o conceito 5, dois eram provenientes de IES Públicas e oito, de IES Privadas.
Dos três cursos participantes de IES Públicas, o conceito 5 foi o valor modal, atribuído a dois cursos ( 66,7 \% ).
O outro curso mantido por IES Pública foi avaliado com conceito 2.
Nesta categoria, nenhum curso ficou sem conceito ( SC ), ou recebeu os conceitos restantes.
Na rede privada, o conceito modal foi 3, com 29 cursos dos 79 da categoria.
Entre os demais cursos participantes, três receberam conceito 1, e 15, conceito 2.
O conceito 4 foi atribuído a 23 cursos e o conceito 5, a oito cursos.
Nesta Categoria Administrativa, apenas um curso ficou sem conceito ( SC ).
Na análise por região, observa -se que, na região Norte, as Instituições Privadas responderam por todos os dois cursos da região, dos quais nenhum ficou sem conceito ( SC ) ou recebeu conceitos 1, 4 ou 5.
Um curso foi alocado no conceito 2 e o outro, no conceito 3.
Não houve registro de nenhum curso oferecido por Instituições Públicas nesta combinação de categoria e região.
Na região Nordeste, a rede privada concentrou 12 dos 14 cursos participantes, o equivalente a 85,7 \% do total da região, proporção que, embora elevada, representou a menor desta categoria entre as regiões do país.
Dentre os cursos oferecidos por IES Privadas no Nordeste, nenhum ficou sem conceito ( SC ) ou foi alocado aos conceitos 1 ou 2.
Sete cursos ( 58,3 \% ) obtiveram conceito 3, o conceito modal, enquanto três cursos receberam conceito 4, e dois cursos ficaram com conceito cinco nesta categoria e região.
As Instituições Públicas dessa região participaram com apenas dois cursos ( 14,3 \% ), e eles se concentraram no conceito 5.
Nenhum curso ficou sem conceito ( SC ) ou recebeu os demais conceitos nesta combinação de categoria e região.
Na região Sudeste, a rede privada respondeu por todos os 38 cursos participantes.
Nesta categoria, na região Sudeste, o conceito modal foi 3 ( 14 cursos ).
Nesta combinação de Categoria Administrativa e Grande Região, apenas um curso não recebeu conceito ( SC ).
Os demais foram avaliados com conceito 1 ( dois cursos ), conceito 2 ( dez cursos ), conceito 4 ( oito cursos ) e conceito 5 ( três cursos ).
Como já mencionado, nenhum curso foi oferecido por Instituições Públicas na região Sudeste.
As Instituições Privadas, também, responderam por todos os 20 cursos participantes da região Sul.
Desses, nove ficaram com conceito 4, o conceito modal.
Nesta combinação de Categoria Administrativa e Grande Região, três cursos ficaram com conceito 2, seis cursos obtiveram conceito 3, enquanto dois cursos receberam conceito 5.
Nenhum curso ficou sem conceito ( SC ) ou recebeu conceito 1.
Assim como nas regiões Norte e Sudeste, as Instituições Públicas na região Sul não ofereceram nenhum curso de Tecnologia em Design de Interiores Na região Centro-Oeste, sete dos oito cursos participantes eram de Instituições Privadas ( 87,5 \% em termos regionais ).
Destes, três concentraram -se no conceito 4, conceito modal.
Nesta combinação de categoria e região, dois cursos foram avaliados com conceitos 1 e 2 ( um curso, \textbf{cada} ).
Quantitativo repetido para os conceitos 3 e 5 ( um curso \textbf{cada} ).
Nenhum curso ficou sem conceito ( SC ) nesta categoria.
Somente um curso foi oferecido por Instituição Pública, e este obteve conceito 2.
Nenhum curso de IES pública ficou sem conceito ( SC ) ou recebeu os demais conceitos dessa categoria e região.
O Gráfico 5.2 apresenta a distribuição do Conceito Enade segundo a Categoria Administrativa da IES.
Os cursos em IES Públicas ( linha azul ) apresentam uma poligonal

% matched lemmas: cada
\end{textsample}
